\documentclass[11pt]{article}
\usepackage[a4paper,margin=25mm]{geometry}
\usepackage{amsmath,amssymb,amsthm,mathtools,hyperref,xurl}
\setlength{\emergencystretch}{2em}
\newtheorem{theorem}{Theorem}
\newtheorem{lemma}[theorem]{Lemma}
\newtheorem{corollary}[theorem]{Corollary}
\newtheorem{proposition}[theorem]{Proposition}
\newcommand{\Z}{\mathbb Z}
\newcommand{\Q}{\mathbb Q}
\title{An exact Pell certificate for a power of two\\and central-binomial divisibility}
\author{\(\mathbb Z\mathbb Q\) Diophantine Project}
\date{23 November 2026}
\begin{document}
\maketitle
\begin{abstract}
We independently prove a polynomial certificate for the conjunction
that \(b\) is a positive power of two and that an integer \(M\)
divides \(\binom{2r}{r}\).  The hypotheses are only
\(r\ge1\), \(M\ge2\), and \(2\le b\le2M\).
The proof identifies two Pell indices exactly, gives a uniform
one-sided approximation with an explicit error bound, and proves
both directions by integer rounding.  The specialization \(M=n^2\)
uses \(p=2ws^2r^2n^6\); this differs from the \(n^2\) printed
in the historical Jones display.  The theorem is a proved local
certificate, not an assertion that the historical display has been
independently proved universal.  Its quadratic implementation costs
no more coordinates than the corresponding previously compiled block.
\end{abstract}

\section{The polynomial certificate}

For an integer \(A\ge2\), define the Pell sequences by
\[
 X_j(A)+Y_j(A)\sqrt{A^2-1}
       =(A+\sqrt{A^2-1})^j\qquad(j\ge0).
\]
Thus \(X_0=1,X_1=A,Y_0=0,Y_1=1\), and each sequence
satisfies \(Z_{j+2}=2AZ_{j+1}-Z_j\).

\begin{theorem}[A generalized binomial certificate]\label{main}
Let \(r,M,b\) be integers satisfying
\[
                    r\ge1,\qquad M\ge2,\qquad 2\le b\le2M.
\]
Then the following are equivalent.
\begin{enumerate}
\item \(b=2^e\) for an integer \(e\ge1\), and
\(M\mid\binom{2r}{r}\).
\item There are positive integers
\[
                 w,s,a,p,k,\tau,h,c,d,\varphi,\gamma,\eta,i,F,j,o
\]
satisfying the following equations:
\begin{align}
 a&=(wM+1)rsM,&p&=2ws^2r^2M^3,\label{ap}\\
 k&=r+1+h(p-1),&\tau^2&=(p^2-1)k^2+1,\label{small}\\
 4(c-ksM)^2+\eta&=k^2,&c&=2r+1+\varphi,\label{close}\\
 d&=bw+(a-2)c+(4a-5)\gamma,\label{linear}\\
 d^2&=(a^2-1)c^2+1,\label{pell1}\\
 F^2&=(a^2-1)i^2c^4+1,\label{pell2}\\
 (d+oF)^2&=(G^2-1)(2r+1+jc)^2+1,\label{pell3}\\
 G&=a+F^2(c+1-a).\label{bridge}
\end{align}
Here \(G\) abbreviates the displayed expression and need not be
regarded as a separate variable.  The alternative expression
\(G=a+F^2(d^2-a)\) gives the same equivalence.
\end{enumerate}
In every positive solution the following identities are forced:
\begin{gather}
 bw=2^{2r+1},\qquad
 c=Y_{2r+1}(a),\qquad d=X_{2r+1}(a),\qquad
 k=Y_{r+1}(p),\label{forced}\\
 sM=\sum_{j=0}^{r}\binom{2r}{j}(wM)^{r-j}.
 \label{forcedY}
\end{gather}
\end{theorem}

The bound \(b\le2M\) allows nonsquare moduli and, for example,
\(b=4,M=2,r=1\).  In the application with \(M=n^2\),
the earlier assumptions \(n\ge b\ge2\) imply the present
hypotheses with room to spare.

\section{The exact Pell index}

We include the elementary index argument so that the proof of
Theorem~\ref{main} does not depend on an unproved source assertion.

\begin{lemma}[Pell classification and elementary estimates]\label{basic}
For \(A\ge2\), every positive integer solution of
\(x^2-(A^2-1)y^2=1\) is
\((X_t(A),Y_t(A))\) for a unique integer \(t\ge1\).
For \(t\ge2\),
\[
 (2A-1)^{t-1}<Y_t(A)\le(2A)^{t-1},\qquad
 Y_t(A)>A^{t-1}.
\]
Also \(Y_t(A)\equiv t\pmod2\), and
\(Y_t(A)\equiv t\pmod L\) whenever \(A\equiv1\pmod L\).
\end{lemma}
\begin{proof}
For classification, if \(y>1\), multiplication by
\(A-\sqrt{A^2-1}\) sends \((x,y)\) to
\[
       (Ax-(A^2-1)y,\ Ay-x).
\]
Both coordinates are positive integers, the norm remains one,
and the second coordinate is smaller than \(y\).  Indeed
\(\sqrt{A^2-1}\,y<x<Ay\), and \(\sqrt{A^2-1}>A-1\).
Descent ends at \((A,1)\).  The recurrence proves strict
monotonicity and uniqueness.  Since \(0<Y_{t-1}<Y_t\) for
\(t\ge2\), the recurrence gives
\((2A-1)Y_t<Y_{t+1}<2AY_t\).  Starting with \(Y_2=2A\)
gives the bounds, including \(Y_t>A^{t-1}\).
Reducing the recurrence modulo two, or modulo \(L\) with
\(A\equiv1\pmod L\), proves the congruences.
\end{proof}

\begin{lemma}[An odd prescribed index]\label{index}
Let \(a\ge2\), let \(B\ge3\) be odd, and let \(c>B,d>0\)
be integers.  Equations \eqref{pell1}--\eqref{pell3}, with
\(2r+1\) replaced by \(B\), have positive integer witnesses
\(i,F,j,o\) if and only if
\[
                         c=Y_B(a),\qquad d=X_B(a).
\]
Either expression for \(G\) in Theorem~\ref{main} is allowed.
\end{lemma}
\begin{proof}
Write \(D=a^2-1\).  By classification,
\[
 c=Y_t(a),\quad d=X_t(a),\qquad ic^2=Y_q(a),\quad F=X_q(a).
\]
We first show that \(tc\mid q\).  Write \(q=vt+u\),
\(0\le u<t\).  Modulo \(c\), the addition formula gives
\(Y_q\equiv d^vY_u\).  The norm gives \(\gcd(c,d)=1\),
and \(0\le Y_u<c\), so \(u=0\).  Expansion of
\((d+c\sqrt D)^v\) then gives
\(Y_q\equiv vd^{v-1}c\pmod{c^2}\).  Thus \(c\mid v\).
In particular \(c\mid q\) and \(q\ge c>t\); here
\(t<c\) follows from \(c>B\ge3\) and Lemma~\ref{basic}.

Both bridges satisfy
\[
             G>a,\qquad G\equiv1\pmod c,
                       \qquad G\equiv a\pmod F.
\]
For the new bridge use \(c\ge Y_2(a)=2a\); for the old
one use \(d^2\equiv1\pmod c\).  The last norm equation
has a positive Pell index \(v\), so
\[
             B+jc=Y_v(G),\qquad d+oF=X_v(G).
\]
It follows that \(v\equiv B\pmod c\) and
\(X_v(a)\equiv X_t(a)\pmod{X_q(a)}\).

The latter congruence implies \(v\equiv\pm t\pmod{4q}\).
For completeness, the doubling identities give
\(X_{2q+u}\equiv-X_u\pmod{X_q}\), while \(X_{-u}=X_u\).
These reduce each residue to \(\pm X_u\), \(0\le u\le q\).
The positive values are distinct in \([1,X_q]\), and the
negative possibility cannot equal \(X_t\), since
\(0<X_u+X_t\le2X_{q-1}<X_q\) when \(u<q\).
The case \(u=q\) is zero modulo \(X_q\) and is also impossible.
This proves the stated index congruence.

Thus \(B\equiv\pm t\pmod c\).  Since \(0<B,t<c\),
either \(B=t\) or \(B=c-t\).  The second alternative is
excluded by \(c\equiv t\pmod2\) and the oddness of \(B\).
This proves necessity.

Conversely suppose \(c=Y_B(a),d=X_B(a)\).  Set \(q=Bc\).
The binomial expansion of \((d+c\sqrt D)^c\) shows that
\(c^2\mid Y_q(a)\).  Set
\[
 i=Y_q(a)/c^2,\quad F=X_q(a),\quad
 j=(Y_B(G)-B)/c,\quad o=(X_B(G)-d)/F.
\]
The displayed congruences for \(G\) show that the last two
quotients are integers.  They are positive because \(G>a\),
\(B\ge3\), \(Y_B(G)>B\), and \(X_B(G)>X_B(a)\).
The norm identities give all three equations.
\end{proof}

\section{A uniform one-sided approximation}

\begin{lemma}[Normalized Pell values]\label{normal}
Put \(A_t(x)=Y_t(x)/(2x)^{t-1}\) for real \(x\ge2\) and
integers \(t\ge1\).  Then
\[
              1-\frac{t-2}{4x^2}\le A_t(x)\le1
                         \qquad(t\ge2).
\]
For fixed \(x\), the sequence is strictly decreasing from
\(t=2\) onward; for fixed \(t\), it is nondecreasing in \(x\).
\end{lemma}
\begin{proof}
The recurrence gives
\[
        A_1=A_2=1,\qquad
        A_{t+1}=A_t-\frac{A_{t-1}}{4x^2}.
\]
The Pell values are positive, so the normalized sequence decreases
and remains positive.  Summing the successive differences and using
\(A_{t-1}\le1\) gives the lower bound.

For monotonicity in \(x\), the recurrence and the identity
\(Y_t(\cos\theta)=\sin(t\theta)/\sin\theta\) show that
the roots of the polynomial \(Y_t(x)\) are
\(\cos(j\pi/t)\), \(1\le j<t\), and its leading
coefficient is \(2^{t-1}\).  Pairing opposite roots gives
\[
 A_t(x)=\prod_{j=1}^{\lfloor(t-1)/2\rfloor}
       \left(1-\frac{\cos^2(j\pi/t)}{x^2}\right).
\]
Every factor is positive and nondecreasing for \(x\ge2\).
The possible zero root contributes the factor one.
\end{proof}

\begin{lemma}[Ratio estimate]\label{ratio}
Let \(r\ge1\), \(Y,V\ge2\) be integers, and put
\[
 a=rY(V+1),\quad p=2r^2Y^2V,\quad
 C=Y_{2r+1}(a),\quad K=Y_{r+1}(p),\quad
 \rho=\frac{(V+1)^{2r}}{V^r}.
\]
Then
\[
        0<\rho-\frac CK<\frac{\rho}{Y^2V^2}.
\]
\end{lemma}
\begin{proof}
Since \(p>a\), Lemma~\ref{normal} gives
\(A_{2r+1}(a)<A_{r+1}(p)\).  The leading terms cancel as
\[
          \frac{(2a)^{2r}}{(2p)^r}=\rho.
\]
This proves the strict upper estimate \(C/K<\rho\).

Put \(\varepsilon=1/(2Y^2V^2)\le1/32\).
The deficits in the two normalized Pell values satisfy
\[
 \frac{2r-1}{4a^2}<\varepsilon,
 \qquad
 \frac{r-1}{4p^2}<\varepsilon.
\]
For the first inequality substitute \(a=rY(V+1)\) and use
\((2r-1)V^2<2r^2(V+1)^2\); for the second substitute
\(p=2r^2Y^2V\) and use \(r-1<8r^4Y^2\).
Thus both normalized values lie in \((1-\varepsilon,1]\).
Their ratio differs from one by less than
\(\varepsilon/(1-\varepsilon)<2\varepsilon\).
Multiplying by \(\rho\) proves the asserted lower estimate.
\end{proof}

\section{Proof of the certificate}

\begin{proof}[Proof of Theorem~\ref{main}: necessity]
Suppose a positive solution is given.  Put
\[
              B=2r+1,\qquad Y=sM,\qquad V=wM.
\]
Then \(Y,V\ge2\), and \eqref{ap} becomes
\(a=rY(V+1)\), \(p=2r^2Y^2V\).
Lemma~\ref{index} gives \(c=Y_B(a),d=X_B(a)\).
The inequality in \eqref{close} says
\begin{equation}\label{half}
                   \left|\frac ck-Y\right|<\frac12.
\end{equation}

We next determine the other Pell index.  From \eqref{small},
\(k=Y_\ell(p)\) for a positive integer \(\ell\), and
\(\ell\equiv r+1\pmod{p-1}\).  The defining formulas give
\[
           p\ge16r^2,\qquad p-1>r+1,\qquad 2a<p^2.
\]
For the last inequality,
\(p^2/(2a)=2r^3Y^3V^2/(V+1)>1\).
If \(\ell\ne r+1\), then
\(\ell\ge p+r\ge4r+1\).  Lemma~\ref{basic} would imply
\[
       k>p^{\ell-1}\ge p^{4r},\qquad
       c<(2a)^{2r}<p^{4r}.
\]
This contradicts \(c>k(Y-1/2)>k\), from \eqref{half}.
Therefore \(\ell=r+1\).

Let \(\rho\) be as in Lemma~\ref{ratio}, and put
\(E=\rho-c/k\).  That lemma gives
\begin{equation}\label{Ebound}
       0<E<\frac{\rho}{Y^2V^2},\qquad
       \delta:=\frac1{Y^2V^2}\le\frac1{16}.
\end{equation}
Since \(\rho=(V+2+1/V)^r>4^r\), equations
\eqref{half}--\eqref{Ebound} imply
\[
       Y>\frac{15}{16}4^r-\frac12
        >\frac{2}{3} 4^r.
\]
The final inequality holds already for \(r=1\): its difference
is \((13/48)4^r-1/2>0\).  Consequently
\begin{equation}\label{a-bounds}
      a=rY(V+1)>2\cdot4^r=2^B,
      \qquad a\ge2(V+1)>2V\ge bw.
\end{equation}
The last inequality uses \(b\le2M\).

Set \(R_t=X_t(a)-(a-2)Y_t(a)\).  The addition identities give
\[
 R_0=1,\quad R_1=2,\qquad
 R_{t+1}=2R_t+(4a-5)Y_t(a).
\]
Hence \(R_B\equiv2^B\pmod{4a-5}\).
Equation \eqref{linear} gives the same residue for \(bw\).
Both \(bw\) and \(2^B\) are strictly between zero and \(a\),
while \(4a-5>a\), since \(a\ge6\).  Their congruence therefore
forces \(bw=2^B\).  In particular \(b\ge2\) is a power of two.

Expand the rational number \(\rho\) as
\[
 \rho=T+\theta,\qquad
 T=\sum_{j=0}^r\binom{2r}{j}V^{r-j},\qquad
 \theta=\sum_{j=r+1}^{2r}\binom{2r}{j}V^{r-j}.
\]
Here \(T\) is an integer, and symmetry of the binomial
coefficients gives
\begin{equation}\label{tail}
 0<\theta<\frac{4^r}{2V}
        =\frac{b}{4M}\le\frac12.
\end{equation}
The inequality is strict even when \(r=1\), since
\(\sum_{j=r+1}^{2r}\binom{2r}{j}<4^r/2\).

Equations \eqref{half}--\eqref{Ebound} also give
\[
 \frac\rho Y<\frac{1+1/(2Y)}{1-\delta}
             \le\frac{5/4}{15/16}=\frac43,
 \qquad E<\frac4{3YV^2}\le\frac16.
\]
Since \(c/k=T+\theta-E\), equation \eqref{half} places
the integer \(Y-T\) strictly between \(-2/3\) and \(1\).
Thus \(Y=T\), proving \eqref{forcedY}.
As \(M\mid Y\) and \(M\mid V\), reduction of this identity
modulo \(M\) gives \(M\mid\binom{2r}{r}\).
\end{proof}

\begin{proof}[Proof of Theorem~\ref{main}: sufficiency]
Suppose \(b=2^e\), \(e\ge1\), and
\(M\mid\binom{2r}{r}\).  Since
\[
      b\le2M\le2\binom{2r}{r}<2\cdot4^r=2^{2r+1},
\]
the integer
\[
                         w=2^{2r+1}/b
\]
is positive.  Put \(V=wM\) and define
\[
 T=\sum_{j=0}^r\binom{2r}{j}V^{r-j},\qquad
 Y=T,\qquad s=T/M.
\]
The divisibility hypothesis and \(M\mid V\) show that \(s\)
is an integer.  It is positive, and \(Y\ge V^r\ge2\).
Define \(a,p\) by \eqref{ap}, and set
\[
 c=Y_{2r+1}(a),\quad d=X_{2r+1}(a),\quad
 k=Y_{r+1}(p),\quad \tau=X_{r+1}(p).
\]
The congruence \(Y_{r+1}(p)\equiv r+1\pmod{p-1}\)
gives an integer \(h=(k-r-1)/(p-1)>0\).
Also \(\varphi=c-(2r+1)>0\).

The same expansion \(\rho=T+\theta\) satisfies
\(0<\theta<1/2\), by \eqref{tail}.  Lemma~\ref{ratio} gives
\[
 0<E=\rho-c/k<\frac{\rho}{Y^2V^2}
       =\frac{1+\theta/Y}{YV^2}
       <\frac5{4YV^2}\le\frac5{32}<\frac12.
\]
Thus \(-1/2<\theta-E=c/k-Y<1/2\), and
\(\eta=k^2-4(c-kY)^2\) is a positive integer.

The recurrence for \(R_t\) gives
\(R_{2r+1}\equiv2^{2r+1}=bw\pmod{4a-5}\), so
\[
                    \gamma=\frac{R_{2r+1}-bw}{4a-5}
\]
is an integer.  It is positive: \(R_2=4a-1>a\), the
sequence \(R_t\) is increasing for \(t\ge1\), and
\(a\ge2(V+1)>2V\ge bw\).  Lemma~\ref{index} supplies
positive integers \(i,F,j,o\) for either bridge.
All the required equations now hold.
\end{proof}

\section{Consequences for the compiled systems}

\begin{corollary}[The square-modulus specialization]
Let \(r\ge1\), \(n\ge b\ge2\).  In Theorem~\ref{main},
put \(M=n^2\).  The resulting polynomial block is equivalent to
\[
              b\text{ is a power of two},\qquad
                     n^2\mid\binom{2r}{r}.
\]
Its auxiliary Pell parameter is
\[
                         p=2ws^2r^2n^6.
\]
\end{corollary}

\begin{proposition}[No additional quadratic coordinates]
Suppose the older quadratic core has already retained \(n^2\)
as a coordinate and implemented the six-coordinate block
\[
 L_0=sr,\quad L_{0,2}=L_0^2,\quad U=wn^2,\quad
 p=2U L_{0,2},\quad Y=sn^2,
\]
\[
                    2L_0a=(p+2L_{0,2})n^2.
\]
After the previously made elimination of the temporary product
\(rY\), the replacement block can use six coordinates
\(L,L_2,U,p,Y,a\) with equations
\begin{gather}
 Y=sn^2,\qquad L=rY,\qquad L_2=L^2,
 \qquad U=wn^2,\label{quad1}\\
                  p=2UL_2,\qquad a=(U+1)L.
                  \label{quad2}
\end{gather}
These are quadratic and imply exactly
\(a=(wn^2+1)rsn^2\), \(p=2ws^2r^2n^6\), for all
integer values, including zero and negative values.
Thus the correction of this local certificate
requires no increase in the six coordinates or six equations.
\end{proposition}
\begin{proof}
The first four equations give \(L=rsn^2\).
The two equations in \eqref{quad2} give the required \(p\) and
\(a\) directly, without division.  Conversely the defining values satisfy
all six equations.  Every displayed equation has total degree
at most two in its coordinates, with \(n^2\) read as the
already available square coordinate.  In the old block the last
equation gives \(a=(U+1)L_0n^2\) because \(L_0=sr>0\).
The six-coordinate comparison refers to the already reduced core,
not an unreduced expression containing a seventh temporary product.
\end{proof}

The positivity reductions proved in the companion manuscript
\emph{An exact three-Pell index certificate and nine fewer quadratic
coordinates} apply to this \(n^6\) block as well as to the earlier
\(n^2\) block.  In particular, the new bridge and the one-square
choice of the positive root do not change Theorem~\ref{main}.

\paragraph{Source status.}
The historical page \cite[Theorem 3, p.~861]{Jones80}, inspected as
an image, prints \(p=2ws^2r^2n^2\).  We have not changed that
transcription or the meaning of the repository's hypothesis
\(\mathcal U_{\rm em}\).  The theorem in this paper is instead
an independently proved replacement of the final local block.
It does not prove that the first, coding part of the historical
system is universal, nor does it verify the alleged 57-variable,
47-quadratic source.  Promoting a complete corrected source to an
undecidability theorem still requires a precise audit of that coding
part and its underlying universality argument.  The giant power
\(q=b^{5^{60}}\) has not been reduced or altered here.

\paragraph{Reproducibility.}
The exact checker
\nolinkurl{undecidability/verify_binomial_pell_certificate.py}
verifies the polynomial identities, the coefficient-neutral
quadratic replacement, finite exact ratio examples, and explicit
binomial examples.  It also expands one full three-Pell witness.
These finite calculations illustrate and audit the formulas; the
proof for all parameters is the argument above.

\begin{thebibliography}{9}
\bibitem{Jones80}
J. P. Jones, \emph{Undecidable Diophantine equations},
Bulletin of the American Mathematical Society (New Series)
\textbf{3}, no.~2 (1980), 859--862.
\url{https://doi.org/10.1090/S0273-0979-1980-14832-6}.
\bibitem{ProjectPell}
\(\mathbb Z\mathbb Q\) Diophantine Project,
\emph{An exact three-Pell index certificate and nine fewer quadratic
coordinates}, 23 November 2026.
\nolinkurl{undecidability/papers/pell_index_optimization.tex}.
\end{thebibliography}
\end{document}
